\clase{7}{25 de Marzo}{}

Se denota la suma directa de $Y$ y $Z$ por $Y \oplus Z$.

\section{II. Funcionales Lineales}

\begin{prop}
	Sean $X$ e.v.n, $f \in X'$. Luego, $f \in X^{*}$ si y solo si $f(B_{1}(0)) \subsetneq \mbb{K}$.
\end{prop}
\begin{proof}[Proof Other Information]
	$\boxed{\Rightarrow}$ $\forall x \in B_{1}(0),\ \| x \| \leq 1 \implies |f(x)| \leq \| f \| \| x \| \leq \| f \|$. Es decir, $f(B_{1}(0)) \subset B_{\| f \|}(0) \subsetneq \mbb{K}$. \par
	$\boxed{\Leftarrow}$ Supongamos que $f$ no es continua. Observar que $\forall \lambda \in \mbb{K},\ |\lambda| \leq 1,\ B_{|\lambda|}(0) = |\lambda| B_{1}(0)$, por lo que $f(|\lambda| B_{1}(0)) = f(\lambda B_{1}(0)) = \lambda f(B_{1}(0)) \subset f(B_{1}(0))$. Luego, existe $(x_{n})_{n \in \N}$ tal que
	\[ \begin{cases}
		\| x_{n} \| = 1 \quad \forall n \in \N \\
		|f(x_{n})| \longrightarrow \infty
	\end{cases} \]
	Luego,
	\[ \xi_{n} = \{\lambda x_{n} \ : \ |\lambda| \leq 1\} \subset B_{1}(0). \]
	Notar que
	\begin{align*}
		f(\xi_{n}) &= \{\lambda f(x_{n}) \ : \ |\lambda| \leq 1\} \\
		&= \Big{\{}\lambda \frac{f(x_{n})}{|f(x_{n})|} |f(x_{n})| \ : \ |\lambda| \leq 1 \Big{\}} \\
		&= B_{|f(x_{n})|}(0)
	.\end{align*}
	Por lo tanto, $f(\xi_{n}) \subset f(B_{1}(0)) \quad \forall n \in \N$. Es decir,
	\[ \bigcup_{n \in \N} f(\xi_{n}) \subset f(B_{1}(0)) \subset \mbb{K} \]
	Además, notar que
	\[ \bigcup_{n \in \N} f(\xi_{n}) = \bigcup_{n \in \N} B_{|f(x_{n})|}(0) = \mbb{K}. \]
	Es decir, $f(B_{1}(0)) = \mbb{K}$.
\end{proof}

\begin{definition}[hiperplano]
	Un hiperplano (afin) de un e.v $X$ es un subconjunto $H \subset X$ de la forma: $H = x + Y$, donde $Y$ es un s.e.v de $X$ de codimensión 1 y $x \in X$.
\end{definition}

\begin{notation}
	Si $f \in X'$, donde $X$ es e.v,
	\[ J(f) \coloneq \{x \in X \ : \ f(x) = 1\} = f^{-1}(\{1\}). \]  
\end{notation}

\begin{prop}
	Sea $X$ un e.v
	\begin{enumerate}[a)]
		\item Si $f \in X',\ f \neq 0$, i.e. $\exists x_{0} \in X,\ f(x_{0}) \neq 0$, entonces $\ker f$ es un s.e.v de $X$ de codimensión 1. Además, $\forall x \in X,\ \exists! (\lambda, y) \in \mbb{K} \times \ker f$ tal que $x = \lambda x_{0} + y$, $J(f)$ es un hiperplano de $X$. 

		\item Si $f,g \in X' \setminus \{0\}$, entonces
		\[ \begin{cases}
			f = \lambda g \\
			\lambda \neq 0
		\end{cases} \iff \ker f = \ker g. \]

		\item El mapeo $f \mapsto J(f)$ es un mapeo 1-1 entre los funcionales de $X' \setminus \{0\}$ y los hiperplanos que no contienen a $0_{X}$. 
	\end{enumerate}
\end{prop}
\begin{proof}[Proof Other Information]
	(a) $\ker f$ s.e.v de $X$. $\checkmark$ Se considera ahora $\hat{f} \colon X \slash \ker f \to \mbb{K}$ tal que $(x + \ker f =)\ [x] \mapsto \hat{f}([x]) = f(x)$. Luego, $\hat{f}$ es lineal e inyectiva (de hecho, es un isomorfismo). Entonces, $\dim_{\mbb{K}}(X \slash \ker f) = 1$. \par
	Para ver la existencia: sean
	\[ x = \underbrace{\frac{f(x)}{f(x_{0})}}_{\lambda} x_{0} + \underbrace{\left( x - \frac{f(x)}{f(x_{0})} x_{0} \right)}_{y}. \]
	y
	\[ f(y) = f \left( x - \frac{f(x)}{f(x_{0})} x_{0} \right) = f(x) - \frac{f(x)}{f(x_{0})} f(x_{0}) = 0. \]
	Es decir, $y \in \ker f$. \par
	Para ver unicidad: supongamos que $x = \lambda x_{0} + y = \lambda' x_{0} + y'$, donde $\lambda, \lambda' \in \mbb{K}$ y $y,y' \in \ker f$. Entonces
	\begin{align*}
		(\lambda - \lambda')x_{0} = y - y' \in \ker f &\implies f((\lambda - \lambda') x_{0}) = 0 = (\lambda - \lambda') f(x_{0}) \\
		&\implies \lambda = \lambda' \\
		&\implies y = y'
	.\end{align*}
	Observamos, además, que
	\begin{align*}
		[x] &= [\lambda x_{0}] \\
		&= \lambda [x_{0}] \\
		&= \lambda x_{0} + \ker f
	\end{align*}
	en la descomposición anterior. Luego,
	\begin{align*}
		J(f) &= \{x \in X \ : \ f(x) = 1\} \\
		&= \{x \in X \ : \ \lambda f(x_{0}) = 1\} \\
		&= \{x \in X \ : \ \hat{f}([x]) = 1\} \\
		&= \lambda_{x} x_{0} + \ker f
	.\end{align*}
	\par (b) $\boxed{\Rightarrow}$ Sea $f = \beta g$, donde $f,g \neq 0$ y $\beta \neq 0$. Luego, $\ker f = \ker g$. \par
	$\boxed{\Leftarrow}$ Supongamos que $\ker f = \ker g$. Por lo visto en a), podemos fijar $x_{0} \not\in \ker f$ tal que $\forall x \in X,\ x = \lambda x_{0} + y$, donde $\lambda \in \mbb{K},\ y \in \ker f$. Notar que $f(x) = \lambda f(x_{0})$. Luego, en
	\[ g(x) = \lambda g(x_{0}) + g(y) \]
	tenemos que $g(y) = 0$, pues $\ker f = \ker g$, y además, $g(x_{0}) \neq 0$ necesariamente (pues $g \neq 0$). Luego,
\end{proof}
